\begin{center}
	\Large{\textbf{EXECUTIVE SUMMARY}}
\end{center}
\vspace{1em}
\normalsize

Searching for medical information online is often a mess of outdated or irrelevant results. I developed this question-answering system to solve this by anchoring its answers in a verified medical knowledge base. By moving away from the model's general training memory and toward specific, curated documents, the system provides more reliable and clinically relevant guidance.

The retrieval stage fuses \ac{bm25} sparse search with dense vector search through \ac{rrf} score combination. Generation is handled by TinyLlama 1.1B, running on \ac{cpu} without any cloud inference or external \ac{api} calls. A five-signal confidence scorer sits on top of the generation step, backed by a DeBERTa \ac{nli} model that flags hallucinated claims before the answer reaches the user.

The knowledge base draws from three sources: ChatDoctor-HealthCareMagic patient-doctor conversations, PubMedQA biomedical research abstracts, and MedMCQA postgraduate medical entrance questions. After processing with a custom RecursiveSentenceChunker, the corpus contains 505,584 documents. The chunking strategy alone accounted for a 0.292 point gain in keyword coverage, pushing the score from 0.3845 under flat character-based chunking to 0.6765 under sentence-boundary-aware chunking, measured on a 97-question evaluation set.

Three generation configurations were tested. The base TinyLlama model scored 0.3845 keyword coverage. \ac{qlora} fine-tuning over 500 steps on MedMCQA and PubMedQA training pairs raised that to 0.4634. BioMistral-7B (GGUF Q4\_K\_M\allowbreak{} quantisation) produced the best ROUGE-L score, 0.299 against TinyLlama's 0.208, but its more discursive output style resulted in lower keyword coverage. On the 36-question end-to-end benchmark, the retrieval pipeline achieved a 97.2\% hit rate.

The full system runs as a FastAPI service with a React frontend and starts with a single \texttt{make\allowbreak{} run} command. No \ac{gpu} is required.
